% Part: intuitionistic-logic
% Chapter: tableaux
% Section: soundness

\documentclass[../../../include/open-logic-section]{subfiles}

\begin{document}

\olfileid{int}{tab}{sou}

\olsection{સ્ફુરણાવાદી \usetoken{P}{tableau}ની યથાર્થતા}

\begin{explain}
  સ્ફુરણાવાદી !!{tableau}s યથાર્થ છે તે બતાવવા,
  જો
  \[
  \sFmla{\True}{!B_1}[1], \dots, \sFmla{\True}{!B_n}[1], \sFmla{\False}{!A}[1]
  \]
  માટે બંધ !!{tableau} હોય, તો
  $!B_1, \dots, !B_n \Entails !A$ છે તે બતાવવું પડે.
  પ્રતિપક્ષી વિધાન સાબિત કરવું સરળ છે: જો કોઈ
  $\mModel{M}$ અને જગત~$w$ માટે બધા $i=1$,
  \dots,~$n$ માટે $\mSat{M}{!B_i}[w]$ હોય,
  પરંતુ $\mSat/{M}{!A}[w]$ હોય, તો કોઈ
  !!{tableau} બંધ થઈ શકે નહીં.
  \sourcecorrection{OLMD-101}{મૂળમાં અહીં
  પ્રતિનિદર્શમાં $!A$ પણ સત્ય છે એમ લખાયું છે;
  $!A$ ખોટું હોવું જોઈએ, તેથી સત્ય-સંજ્ઞામાં
  નિષેધચિહ્ન ઉમેર્યું છે.}
  આવું પ્રતિનિદર્શ બતાવે છે કે !!{tableau}ની
  શરૂઆતની ધારણાઓ સંતોષ્ય છે. સાબિતીની રીત એ
  બતાવવાની છે કે !!a{tableau}ની શાખા પરનાં બધાં
  ઉપસર્ગવાળા !!{formula}s સંતોષ્ય હોય, તો કોઈ પણ
  નિયમ લાગુ કરતાં ઓછામાં ઓછી એક વિસ્તરેલી શાખા
  સંતોષ્ય રહે છે. બંધ શાખાઓ અસંતોષ્ય હોવાથી,
  સંતોષ્ય ઉપસર્ગવાળા !!{formula}sના ગણનું કોઈ પણ
  !!{tableau} ઓછામાં ઓછી એક ખુલ્લી શાખા ધરાવે છે.

  સ્ફુરણાવાદી કિસ્સામાં આ રીત લાગુ કરવા
  ``સંતોષ્ય''ની વ્યાખ્યાને સંબંધાત્મક નિદર્શો અને
  ઉપસર્ગો સુધી વિસ્તૃત કરવી પડે છે. પછી સાબિતી
  સીધી છે.
  \sourcecorrection{OLMD-102}{મૂળમાં અહીં
  ``મોડલ કિસ્સો'' લખાયો છે; આ પ્રકરણ
  સ્ફુરણાવાદી ટેબ્લોનું છે, તેથી વિષયનું નામ
  સુસંગત કર્યું છે.}
\end{explain}

\begin{defn}
  ધારો કે $P$ ઉપસર્ગોનો ગણ છે, એટલે
  $P \subseteq (\PosInt)^* \setminus \{\emptyseq\}$,
  અને $\mModel{M}$ નિદર્શ છે.
  વિધેય~$f\colon P \to W$ એ~$\mModel{M}$માં
  $P$નું \emph{અર્થઘટન} ત્યારે કહેવાય જ્યારે
  $\sigma$ અને $\sigma.n$ બંને~$P$માં હોય ત્યારે
  $Rf(\sigma)f(\sigma.n)$ હોય.

  ઉપસર્ગો~$P$ના અર્થઘટનની દૃષ્ટિએ નીચેની
  વ્યાખ્યાઓ આપી શકીએ:
  \begin{enumerate}
  \item $\mModel{M}$ એ $\sFmla{\True}{!A}[\sigma]$ને
    ત્યારે અને માત્ર ત્યારે જ સંતોષે જ્યારે
    $\mSat{M}{!A}[f(\sigma)]$.
  \item $\mModel{M}$ એ $\sFmla{\False}{!A}[\sigma]$ને
    ત્યારે અને માત્ર ત્યારે જ સંતોષે જ્યારે
    $\mSat/{M}{!A}[f(\sigma)]$.
  \end{enumerate}
\end{defn}

$R$ સ્વવાચક અને પરંપરિત હોવાથી, અને $\sigma.{*}$
એ $\sigma$, $\sigma.n_1$, $\sigma.n_1.n_2$,
\dots, દર્શાવતું હોવાથી, $Rf(\sigma)f(\sigma.{*})$
પણ છે.

\begin{defn}
  ધારો કે $\Gamma$ ઉપસર્ગવાળા !!{formula}sનો ગણ છે,
  અને $P(\Gamma)$ તેમાં આવતા ઉપસર્ગોનો ગણ છે.
  જો $f$ એ $\mModel{M}$માં $P(\Gamma)$નું અર્થઘટન
  હોય, તો $\mModel{M}$ એ~$f$ની દૃષ્ટિએ $\Gamma$ને
  સંતોષે છે, સંજ્ઞામાં $\mSat{M}{\Gamma}[f]$,
  એમ ત્યારે કહીએ જ્યારે $\mModel{M}$ એ~$f$ની
  દૃષ્ટિએ $\Gamma$નું દરેક ઉપસર્ગવાળું !!{formula}
  સંતોષે. $\Gamma$ \emph{સંતોષ્ય} ત્યારે અને માત્ર
  ત્યારે જ જ્યારે એવું નિદર્શ~$\mModel{M}$ અને
  $P(\Gamma)$નું અર્થઘટન~$f$ હોય કે
  $\mSat{M}{\Gamma}[f]$.
\end{defn}

\begin{prop}
  જો $\Gamma$માં કોઈ !!{formula}~$!A$ અને
  ઉપસર્ગ~$\sigma$ માટે
  $\sFmla{\True}{!A}[\sigma]$ તથા
  $\sFmla{\False}{!A}[\sigma.{*}]$ બંને હોય,
  અથવા તેમાં $\sFmla{\True}{\lfalse}[\sigma]$
  હોય, તો $\Gamma$ અસંતોષ્ય છે.
\end{prop}

\begin{proof}
  હંમેશા $\mSat/{M}{\lfalse}[f(\sigma)]$ હોવાથી,
  $\sFmla{\True}{\lfalse}[\sigma]$ ધરાવતો
  $\Gamma$ અસંતોષ્ય છે.

  વળી, બંને વિરોધી ચિહ્નિત સૂત્રોને સંતોષતું
  નિદર્શ~$\mModel{M}$ અને $P(\Gamma)$નું અર્થઘટન~$f$
  હોઈ શકે નહીં. જો $\mSat{M}{!A}[f(\sigma)]$
  હોય, તો $Rf(\sigma)f(\sigma.{*})$ અને
  \olref[sem][rel]{prop:true-monotonic} મુજબ
  $\mSat{M}{!A}[f(\sigma.{*})]$ પણ હોય.
  તેથી $\mSat{M}{!A}[f(\sigma)]$ અને
  $\mSat/{M}{!A}[f(\sigma.{*})]$ બંને સાથે
  હોઈ શકે નહીં.
  \sourcecorrection{OLMD-103}{મૂળના આ વાક્યમાં
  વાક્યરચના અધૂરી છે, $R$ની બીજી દલીલમાં
  $f$ ખૂટે છે અને એકદિશવર્ધિતાથી મળતું સત્ય
  ફરી મૂળ જગતમાં મૂકાયું છે. અહીં તે સુલભ
  જગત~$f(\sigma.{*})$માં મૂક્યું છે; અસત્યના
  ચિહ્નિત સૂત્રને પણ~$\sigma$ ઉપસર્ગ આપ્યો છે.}
\end{proof}

\begin{thm}[યથાર્થતા]
  \ollabel{thm:tableau-soundness}
  જો $\Gamma$ માટે બંધ !!{tableau} હોય,
  તો $\Gamma$ અસંતોષ્ય છે.
\end{thm}

\begin{proof}
!!a{tableau}ની શાખાને ત્યારે સંતોષ્ય કહીએ જ્યારે
તેના પરના !!{signed formula}sનો ગણ સંતોષ્ય હોય;
અને !!a{tableau}ને ત્યારે સંતોષ્ય કહીએ જ્યારે
તેની ઓછામાં ઓછી એક શાખા સંતોષ્ય હોય.

આ બતાવીએ: સંતોષ્ય !!{tableau}ને અનુમાનનો કોઈ પણ
નિયમ વાપરીને વિસ્તારો, તો મળતું !!{tableau} પણ
સંતોષ્ય રહે છે. આનાથી પ્રમેય મળશે: કોઈ પણ બંધ
!!{tableau} એ~$\Gamma$ની ધારણાઓથી જ શરૂ થતા
!!{tableau} પર અનુમાનના નિયમો લગાવવાથી મળે છે.
જો $\Gamma$ સંતોષ્ય હોય, તો તેના માટેનું દરેક
!!{tableau} સંતોષ્ય રહે. પરંતુ બંધ !!{tableau}
સંતોષ્ય નથી, કારણ કે તેની બધી શાખાઓ બંધ છે
અને બંધ શાખાઓ અસંતોષ્ય છે.

ધારો કે આપણી પાસે સંતોષ્ય !!{tableau} છે,
એટલે ઓછામાં ઓછી એક સંતોષ્ય શાખાવાળું
!!a{tableau}. નિયમ લાગુ કરતાં કોઈ શાખા પર
!!{signed formula}s ઉમેરાય છે અથવા શાખા
બે ભાગમાં ફાટે છે. જો !!{tableau}માં નિયમથી ન વિસ્તરતી
બીજી કોઈ સંતોષ્ય શાખા હોય, તો તે વિસ્તરેલા
!!{tableau}માં પણ સંતોષ્ય રહે છે; તેથી આખું
વિસ્તરેલું ટેબ્લો સંતોષ્ય છે. હવે ફક્ત એવો
કિસ્સો જોવો છે જેમાં નિયમ સંતોષ્ય શાખા પર
લાગુ થાય છે.

તે શાખા પરનાં !!{signed formula}sનો ગણ
$\Gamma$ હોય અને નિયમ જેના પર લાગુ થાય છે તે
!!{signed formula} $\sFmla{S}{!A}[\sigma] \in \Gamma$
હોય. નિયમથી શાખા ન ફાટે, તો વિસ્તરેલી શાખા,
એટલે નિયમના નિષ્કર્ષો સાથેનો $\Gamma$,
હજી સંતોષ્ય છે તે બતાવવું પડે. શાખા ફાટે,
તો બેમાંથી ઓછામાં ઓછી એક નવી શાખા સંતોષ્ય
છે તે બતાવવું પડે.
\begin{enumerate}
\item $\sFmla{\True}{\lnot !B}[\sigma] \in \Gamma$
  પર $\TRule{\True}{\lnot}$ લગાવવાથી શાખા
  વિસ્તરે છે. વિસ્તરેલી શાખામાં
  !!{signed formula}sનો ગણ $\Gamma \cup
  \{\sFmla{\False}{!B}[\sigma.{*}]\}$ હોય છે.
  ધારો કે $\mSat{M}{\Gamma}[f]$. ખાસ કરીને
  $\mSat{M}{\lnot !B}[f(\sigma)]$. તેથી
  $Rf(\sigma)w$ ધરાવતા કોઈ પણ $w$ માટે
  $\mSat/{M}{!B}[w]$; તેમાં $f(\sigma.{*})$
  પણ આવે છે. આથી $\mModel{M}$ એ~$f$ની
  દૃષ્ટિએ $\sFmla{\False}{!B}[\sigma.{*}]$
  સંતોષે છે.
\item $\sFmla{\False}{\lnot !B}[\sigma] \in \Gamma$
  પર $\TRule{\False}{\lnot}$ લગાવવાથી શાખા
  વિસ્તરે છે: સ્વાધ્યાય.
\item $\sFmla{\True}{!B \land !C}[\sigma] \in \Gamma$
  પર $\TRule{\True}{\land}$ લગાવવાથી શાખા
  વિસ્તરે છે અને તે જ શાખા પર બે નવાં
  !!{signed formula}s, $\sFmla{\True}{!B}[\sigma]$
  અને $\sFmla{\True}{!C}[\sigma]$, મળે છે.
  ધારો કે $\mSat{M}{\Gamma}[f]$; ખાસ કરીને
  $\mSat{M}{!B \land !C}[f(\sigma)]$.
  ત્યારે $\mSat{M}{!B}[f(\sigma)]$ અને
  $\mSat{M}{!C}[f(\sigma)]$. એટલે
  $\mModel{M}$ એ~$f$ની દૃષ્ટિએ
  $\sFmla{\True}{!B}[\sigma]$ અને
  $\sFmla{\True}{!C}[\sigma]$ બંને સંતોષે છે.
\item $\sFmla{\False}{!B \lor !C}[\sigma] \in \Gamma$
  પર $\TRule{\False}{\lor}$ લગાવવાથી શાખા
  વિસ્તરે છે: સ્વાધ્યાય.
  \sourcecorrection{OLMD-104}{મૂળની આ
  સ્વાધ્યાય-પંક્તિમાં ચિહ્નિત વિયોજન પછી
  ઉપસર્ગ~$\sigma$ ખૂટે છે; નિયમનો વિષય
  નિશ્ચિત કરવા તે ઉમેર્યો છે.}
\item $\sFmla{\False}{!B \lif !C}[\sigma] \in \Gamma$
  પર $\TRule{\False}{\lif}$ લગાવવાથી શાખા
  વિસ્તરે છે. એ જ શાખા પર બે નવાં
  !!{signed formula}s મળે છે:
  $\sFmla{\True}{!B}[\sigma.n]$ અને
  $\sFmla{\False}{!C}[\sigma.n]$; અહીં
  $\sigma.n$ શાખા પર નવો ઉપસર્ગ છે, એટલે
  $\sigma.n \notin P(\Gamma)$.

  $\Gamma$ સંતોષ્ય હોવાથી એવું
  $\mModel{M}$ અને $P(\Gamma)$નું અર્થઘટન~$f$
  છે કે $\mSat{M}{\Gamma}[f]$; ખાસ કરીને
  $\mSat/{M}{!B \lif !C}[f(\sigma)]$.
  $\Gamma \cup \{\sFmla{\True}{!B}[\sigma.n],
  \sFmla{\False}{!C}[\sigma.n]\}$ સંતોષ્ય છે
  તે બતાવવાનું છે. તે માટે $P(\Gamma)
  \cup \{\sigma.n\}$નું અર્થઘટન નીચે મુજબ
  વ્યાખ્યાયિત કરીએ છીએ:
  \sourcecorrection{OLMD-105}{મૂળ આ લક્ષ્યમાં
  બે નવા ચિહ્નિત સૂત્રોને બદલે
  $\sFmla{\False}{!B \lif !C}[\sigma.n]$
  લખે છે; એ શાખામાં લાગુ પડેલા નિયમનું
  નિષ્કર્ષ નથી. બંને જરૂરી સૂત્રો મૂક્યા છે.}

  $\mSat/{M}{!B \lif !C}[f(\sigma)]$
  હોવાથી એવું $w \in W$ છે કે $Rf(\sigma)w$
  અને $\mSat{M}{!B}[w]$ તથા
  $\mSat/{M}{!C}[w]$. $f'$ એ $f$ જેવું જ હોય,
  ફક્ત $f'(\sigma.n) = w$ મૂકો.
  $f'(\sigma) = f(\sigma)$ અને $Rf(\sigma)w$
  હોવાથી $Rf'(\sigma)f'(\sigma.n)$; તેથી
  $f'$ એ $P(\Gamma) \cup \{\sigma.n\}$નું
  અર્થઘટન છે. સ્પષ્ટ છે કે
  $\mSat{M}{!B}[f'(\sigma.n)]$ અને
  $\mSat/{M}{!C}[f'(\sigma.n)]$.
  $\sigma' \in P(\Gamma)$ એવા દરેક ઉપસર્ગ માટે
  $f(\sigma') = f'(\sigma')$ હોવાથી
  $\mSat{M}{\Gamma}[f']$. તેથી $\mModel{M}, f'$
  એ $\Gamma \cup \{\sFmla{\True}{!B}[\sigma.n],
  \sFmla{\False}{!C}[\sigma.n]\}$ સંતોષે છે.
\end{enumerate}
હવે બે નવી શાખાઓ આપતા નિયમો જોઈએ.
\begin{enumerate}
\item $\sFmla{\False}{!B \land !C}[\sigma] \in \Gamma$
  પર $\TRule{\False}{\land}$ લગાવવાથી
  શાખા બેમાં ફાટે છે: ડાબી શાખા
  $\sFmla{\False}{!B}[\sigma]$થી આગળ વધે છે
  અને જમણી $\sFmla{\False}{!C}[\sigma]$થી.
  ધારો કે $\mSat{M}{\Gamma}[f]$; ખાસ કરીને
  $\mSat/{M}{!B \land !C}[f(\sigma)]$.
  ત્યારે $\mSat/{M}{!B}[f(\sigma)]$ અથવા
  $\mSat/{M}{!C}[f(\sigma)]$. પહેલા કિસ્સામાં
  $\mModel{M}, f$ એ $\sFmla{\False}{!B}[\sigma]$
  સંતોષે છે, એટલે ડાબી શાખા સંતોષ્ય છે.
  બીજા કિસ્સામાં $\mModel{M}, f$ એ
  $\sFmla{\False}{!C}[\sigma]$ સંતોષે છે,
  એટલે જમણી શાખા સંતોષ્ય છે.
\item $\sFmla{\True}{!B \lor !C}[\sigma] \in \Gamma$
  પર $\TRule{\True}{\lor}$ લગાવવાથી શાખા
  વિસ્તરે છે: સ્વાધ્યાય.
\item $\sFmla{\True}{!B \lif !C}[\sigma] \in \Gamma$
  પર $\TRule{\True}{\lif}$ લગાવવાથી શાખા
  વિસ્તરે છે: સ્વાધ્યાય.
\end{enumerate}
\end{proof}

\begin{prob}
\olref[int][tab][sou]{thm:tableau-soundness}ની
સાબિતી પૂરી કરો.
\end{prob}

\begin{cor}
\ollabel{cor:entailment-soundness}
જો $\Gamma \Proves !A$ હોય, તો
$\Gamma \Entails !A$.
\end{cor}

\begin{proof}
  જો $\Gamma \Proves !A$ હોય, તો કોઈ
  $!B_1$, \dots, $!B_n \in \Gamma$ માટે
  $\Delta = \{\sFmla{\False}{!A}[1],
  \sFmla{\True}{!B_1}[1], \dots,
  \sFmla{\True}{!B_n}[1]\}$નું બંધ
  !!{tableau} છે. $\Gamma \Entails !A$
  બતાવવું છે. ઊલટું ધારો: કોઈ
  $\mModel{M}$ અને $w$ માટે $i=1$,
  \dots,~$n$ માટે $\mSat{M}{!B_i}[w]$
  હોય, પણ $\mSat/{M}{!A}[w]$.
  $f(1) = w$ મૂકો; ત્યારે $f$ એ
  $P(\Delta)$નું~$\mModel{M}$માં અર્થઘટન
  છે અને $\mModel{M}$ એ~$f$ની દૃષ્ટિએ
  $\Delta$ને સંતોષે છે. પરંતુ
  \olref{thm:tableau-soundness} મુજબ બંધ
  !!{tableau} ધરાવતો $\Delta$ અસંતોષ્ય છે;
  આ વિરોધાભાસ છે. તેથી
  $\Gamma \Entails !A$ હોવું જ જોઈએ.
  \sourcecorrection{OLMD-106}{મૂળની અંતિમ
  પંક્તિમાં ફરી $\Gamma \Proves !A$ લખાયું છે,
  જે પહેલેથી ધારેલું હતું. વિરોધાભાસથી
  અહીં જોઈતું $\Gamma \Entails !A$ મળે છે.}
\end{proof}

\begin{cor}
\ollabel{cor:weak-soundness}
જો $\Proves !A$ હોય, તો બધા નિદર્શોમાં
$!A$ સત્ય છે.
\end{cor}

\end{document}
